This research project demonstrates the feasibility of generating synthetic patient data through rule-based approaches using Synthea and generative AI approaches using LLMs. Synthea-generated records showed significant improvements over LLM generated records concerning longitudinal consistency and a positive trend over real-world patient data. By sequentially importing these records into a HIS using HL7 FHIR standards, the system simulates the continuous nature of real-world clinical documentation. Additionally, continuous data streams from wearables were imported into the HIS and displayed using the patient's native dashboard. These dynamic structures allow students to analyze realistic disease progression and patient-generated data without privacy risks.

Despite these advantages, the current implementation has several limitations. Rule-based profiles from Synthea produced repetitive patterns, and their FHIR bundles were rejected by the HIS due to structural complexity. Conversely, Large Language Models generated simpler, compatible FHIR data but suffered from broken timelines and clinical hallucinations. Furthermore, the sensor integration lacked interactive frontend features like temporal scaling, and the evaluation was limited by low inter-rater reliability.

Future research should focus on technical and methodological refinements. Key priorities include developing scalable translation scripts to map complex rule-based outputs to HIS-compliant schemas, alongside implementing knowledge-grounded AI architectures, such as Retrieval-Augmented Generation and multi-agent frameworks, to actively enforce guideline compliance and eliminate hallucinations. Additionally, dedicated frontend engineering is required to enable dynamic downsampling and event marking within the sensor dashboard. Addressing these technical barriers and expanding future evaluations to full multi-reviewer validation will help establish reliable, dynamic simulation environments for medical education.